% !TEX root = ../../main.tex

\section{Do-Notation}
\label{sec:do-notation}

La \textit{do-notation} es una sintaxis de \textsf{Haskell} para escribir cómputos monádicos de forma secuencial. Un bloque \texttt{do} contiene statements que pueden vincular resultados intermedios y utilizarlos en los statements posteriores. Esta notación no define una forma nueva de composición, sino que proporciona una escritura más directa para expresiones construidas mediante las operaciones monádicas como puede ser el operador \textit{bind} \verb|(>>=)|, \verb|(>>)| y \texttt{fail}, además de contemplar una regla propia para declaraciones \texttt{let} \cite[sección~3.14]{Haskell2010}.

Para observar la motivación de esta sintaxis, se retomará la suma de dos valores de tipo \texttt{Maybe} presentada mediante el operador \textit{bind} en el Código~\ref{lst:monad-maybe-sum}. Aunque la expresión es compacta, la acumulación de operadores, funciones anónimas y paréntesis dificulta distinguir su estructura a simple vista.

El Código~\ref{lst:monad-maybe-sum-multiline} distribuye la misma expresión en varias líneas, sin cambiar el orden ni la agrupación de sus operaciones, para mostrar con mayor claridad el flujo de definición y composición de los cómputos:

\begin{lstlisting}[style=haskell,caption={Forma multilínea de la suma monádica sobre \texttt{Maybe}.},label={lst:monad-maybe-sum-multiline}]
    Just 3 >>= (\x ->
        Just 2 >>= (\y ->
            return ((+) x y)))
\end{lstlisting}

La indentación permite reconocer que la composición con \texttt{Just 2} se encuentra dentro de la función que introduce \texttt{x}, mientras que \texttt{return} cierra la composición utilizando los identificadores \texttt{x} e \texttt{y}. Este cambio es únicamente visual: la expresión conserva la semántica del Código~\ref{lst:monad-maybe-sum} y produce \texttt{Just 5}.

La \textit{do-notation} permite expresar este mismo flujo sin escribir de forma explícita las funciones anónimas ni las aplicaciones sucesivas de \verb|(>>=)|. En un statement de la forma \verb|p <- e|, la expresión \texttt{e} representa el cómputo y el patrón \texttt{p} vincula su resultado para que pueda utilizarse en los statements posteriores. El Código~\ref{lst:do-notation-maybe-sum} presenta la escritura equivalente:

\noindent\begin{minipage}{\textwidth}
\begin{lstlisting}[style=haskell,caption={Suma monádica sobre \texttt{Maybe} escrita mediante \textit{do-notation}.},label={lst:do-notation-maybe-sum}]
    do
        x <- Just 3
        y <- Just 2
        return ((+) x y)
\end{lstlisting}
\end{minipage}

En el Código~\ref{lst:do-notation-maybe-sum}, el primer statement corresponde al cómputo entregado al \verb|(>>=)| exterior y vincula su resultado con \texttt{x}; el segundo corresponde al \verb|(>>=)| interior y vincula su resultado con \texttt{y}; el último construye el resultado del bloque. Al aplicar las reglas convencionales de traducción de la \textit{do-notation}, este bloque recupera la composición anidada del Código~\ref{lst:monad-maybe-sum-multiline} y, en última instancia, la expresión compacta del Código~\ref{lst:monad-maybe-sum} \cite[sección~3.14]{Haskell2010}.

Aunque el bloque se escribe como una secuencia, en este ejemplo la expresión \texttt{Just 2} no utiliza el valor de \texttt{x}; por ello, entre los dos primeros statements no existe una dependencia de datos local. Ambos resultados sí son utilizados por el statement final. Esta distinción entre orden sintáctico y dependencias de datos será relevante más adelante: la notación hace más legible la composición, pero no autoriza por sí sola a intercambiar cómputos. Las condiciones que permiten transformar o reordenar partes de esta secuencia se estudiarán al introducir las mónadas conmutativas y \texttt{ApplicativeDo}.
